\chapter{Quiz}
\label{cap:quiz}

\begin{nota}
Domande in inglese, come all'esame; spiegazioni in italiano. Le domande sono raggruppate per argomento
e ognuna rimanda alla lezione in cui il tema è stato trattato (link cliccabile). Le domande su ENISA,
NIS2 e DORA (strategia) ricostruiscono le domande chiuse dell'appello precedente
(Capitolo~\ref{cap:tema-precedente}); le opzioni sono nostre. Tutte le domande sono scritte dai
redattori, non dal docente. Soluzioni in fondo al capitolo.
\end{nota}

% =====================================================================
\section{Fondamenti}

\begin{domanda}[Fonti del diritto UE -- \rif{2026-09-21}{21/09}]\label{q:regolamento}
An EU \emph{regulation}, unlike a directive:
\begin{opzioni}
  \item must be transposed into national law before it applies;
  \item is directly applicable in all Member States;
  \item binds only EU institutions;
  \item is a form of soft law.
\end{opzioni}
\end{domanda}

\begin{domanda}[Hard e soft law -- \rif{2026-11-09}{09/11}]\label{q:softlaw}
Which of the following is \textbf{soft law}?
\begin{opzioni}
  \item the Cyber Resilience Act;
  \item the GDPR;
  \item ISO/IEC 27001;
  \item the NIS2 Directive as transposed by D.lgs.~138/2024.
\end{opzioni}
\end{domanda}

\begin{domanda}[Acronimi -- \rif{2026-09-21}{21/09}]\label{q:dpa}
In a contract between a controller and a cloud provider, the acronym ``DPA'' most likely refers to:
\begin{opzioni}
  \item the Data Protection Authority;
  \item the Data Processing Agreement under Art.~28 GDPR;
  \item the Data Protection Act of the provider's country;
  \item the Digital Product Assessment under the CRA.
\end{opzioni}
\end{domanda}

\begin{domanda}[Dati personali -- \rif{2026-09-21}{21/09}]\label{q:anonimi}
Which statement about anonymisation is \textbf{correct}?
\begin{opzioni}
  \item Removing name and surname is enough to make data anonymous.
  \item Pseudonymised data are no longer personal data.
  \item Data that are truly anonymous fall outside the scope of the GDPR.
  \item Anonymisation is a legal basis under Art.~6 GDPR.
\end{opzioni}
\end{domanda}

% =====================================================================
\section{GDPR}

\begin{domanda}[GDPR, basi giuridiche -- \rif{2026-09-25}{25/09}]\label{q:gdpr-basi}
Legal bases for processing under Art.~6 GDPR. \textbf{Select the WRONG statement.}
\begin{opzioni}
  \item Consent must be freely given, specific, informed and unambiguous.
  \item Performance of a contract covers only the processing strictly necessary for that contract.
  \item Public authorities may rely on legitimate interest for processing carried out in the
    performance of their tasks.
  \item Different legal bases may apply to different purposes within the same service.
\end{opzioni}
\end{domanda}

\begin{domanda}[GDPR, ruoli -- \rif{2026-09-25}{25/09}]\label{q:ruoli-fatto}
Two companies sign an agreement stating that each is an independent controller, but in practice they
jointly decide why and how the data are processed. Under the GDPR they are:
\begin{opzioni}
  \item independent controllers, because the contract says so;
  \item joint controllers, because the factual situation prevails over the contract;
  \item controller and processor, depending on who pays;
  \item outside the GDPR, since the roles are disputed.
\end{opzioni}
\end{domanda}

\begin{domanda}[GDPR, ruoli -- \rif{2026-11-09}{09/11}]\label{q:autorizzato}
An employee of the controller who processes personal data following the controller's instructions is:
\begin{opzioni}
  \item a processor (Art.~28);
  \item a sub-processor;
  \item a person authorised to process the data under the authority of the controller;
  \item a joint controller.
\end{opzioni}
\end{domanda}

\begin{domanda}[GDPR, legittimo interesse -- \rif{2026-09-28}{28/09}]\label{q:legittimo}
About legitimate interest (Art.~6(1)(f)). \textbf{Select the WRONG statement.}
\begin{opzioni}
  \item It requires a balancing test against the data subject's interests and fundamental rights.
  \item The legitimate interest invoked is the one of the data subject.
  \item It can be relied upon, for example, for some forms of video surveillance.
  \item The data subject has the right to object to processing based on it.
\end{opzioni}
\end{domanda}

\begin{domanda}[GDPR, principi -- \rif{2026-09-25}{25/09}]\label{q:principi}
Which of the following is \textbf{not} a principle listed in Art.~5 GDPR?
\begin{opzioni}
  \item purpose limitation;
  \item data minimisation;
  \item storage limitation;
  \item freedom of contract.
\end{opzioni}
\end{domanda}

\begin{domanda}[GDPR, ambito -- \rif{2026-09-25}{25/09}]\label{q:domestica}
A private individual publishes on a public website information about named colleagues of a
volunteer group. The household exemption:
\begin{opzioni}
  \item applies, because the individual acts for non-commercial purposes;
  \item does not apply, because the data are made accessible to an indefinite number of people;
  \item applies only if the data are not sensitive;
  \item applies only outside working hours.
\end{opzioni}
\end{domanda}

\begin{domanda}[GDPR, sanzioni -- \rif{2026-09-28}{28/09}]\label{q:sanzioni}
For infringements of the security obligations (including data breach notification duties), the GDPR
administrative fines can reach:
\begin{opzioni}
  \item EUR 20 million or 4\% of total worldwide annual turnover;
  \item EUR 10 million or 2\% of total worldwide annual turnover;
  \item EUR 35 million or 7\% of total worldwide annual turnover;
  \item EUR 15 million or 2.5\% of total worldwide annual turnover.
\end{opzioni}
\end{domanda}

\begin{domanda}[GDPR, categorie di dati -- \rif{2026-10-12}{12/10}]\label{q:malattia}
Sick-leave certificates of employees are:
\begin{opzioni}
  \item common personal data, because they do not state the diagnosis;
  \item data concerning health, a special category under Art.~9;
  \item judicial data under Art.~10;
  \item not personal data, because they are produced by a doctor.
\end{opzioni}
\end{domanda}

\begin{domanda}[GDPR, DPIA -- \rif{2026-10-05}{05/10}]\label{q:dpia}
After a DPIA, the residual risk of a new processing remains high despite the planned measures. The
controller must:
\begin{opzioni}
  \item start the processing and record the decision;
  \item consult the supervisory authority before starting the processing (Art.~36);
  \item notify a data breach;
  \item ask the data subjects for their consent, which removes the risk.
\end{opzioni}
\end{domanda}

% =====================================================================
\section{Data breach e severità ENISA}

\begin{domanda}[Data breach, definizione -- \rif{2026-10-05}{05/10}]\label{q:def-breach}
Which of the following is \textbf{not} a personal data breach?
\begin{opzioni}
  \item an unencrypted laptop with customer data is lost on a train;
  \item ransomware makes the HR database unavailable for a day;
  \item a phishing email reaches an employee, who deletes it without clicking;
  \item an employee browses customer records without any business need.
\end{opzioni}
\end{domanda}

\begin{domanda}[Data breach, riservatezza -- \rif{2026-10-19}{19/10}]\label{q:accesso}
An attacker accessed a file server containing personal data; logs show no evidence of exfiltration.
Under the GDPR:
\begin{opzioni}
  \item there is no breach until exfiltration is proven;
  \item there is a confidentiality breach, because unauthorised access is enough;
  \item there is only an availability breach;
  \item the event is relevant only under NIS2.
\end{opzioni}
\end{domanda}

\begin{domanda}[Data breach, 72 ore -- \rif{2026-10-05}{05/10}]\label{q:decorrenza}
The 72-hour deadline of Art.~33 GDPR runs from:
\begin{opzioni}
  \item the moment the breach occurred;
  \item the moment the controller became aware of the breach;
  \item the end of the forensic investigation;
  \item the first business day after the breach.
\end{opzioni}
\end{domanda}

\begin{domanda}[Data breach, responsabile -- \rif{2026-10-02}{02/10}]\label{q:responsabile}
A cloud provider acting as processor suffers a personal data breach. Who notifies the supervisory
authority?
\begin{opzioni}
  \item The processor, directly and instead of the controller.
  \item The controller, after the processor has informed it without undue delay.
  \item Both, independently and with the same deadline.
  \item ENISA, through the Single Reporting Platform.
\end{opzioni}
\end{domanda}

\begin{domanda}[Data breach, registro -- \rif{2026-10-05}{05/10}]\label{q:registro}
About the internal breach register (Art.~33(5)). \textbf{Select the WRONG statement.}
\begin{opzioni}
  \item It must document the facts, the effects and the remedial action taken.
  \item Breaches assessed as low risk and not notified need not be recorded.
  \item It allows the supervisory authority to verify compliance with Art.~33.
  \item A temporary loss of availability is recorded even if no notification is required.
\end{opzioni}
\end{domanda}

\begin{domanda}[Data breach, disponibilità -- \rif{2026-10-12}{12/10}]\label{q:disponibilita}
A ransomware attack makes personal data unavailable for six hours; the data are fully restored from
backup and there is no exfiltration. Under the GDPR this is:
\begin{opzioni}
  \item not a personal data breach, because the data were restored;
  \item a personal data breach (loss of availability), to be at least recorded in the breach register;
  \item a breach only if the ransom was paid;
  \item a NIS2 matter only.
\end{opzioni}
\end{domanda}

\begin{domanda}[Severità ENISA -- \rif{2026-10-02}{02/10}]\label{q:formula}
In the ENISA methodology the severity of a personal data breach is computed as:
\begin{opzioni}
  \item $SE = DPC + EI + CB$
  \item $SE = DPC \times EI + CB$
  \item $SE = (DPC + CB) \times EI$
  \item $SE = \max(DPC, EI, CB)$
\end{opzioni}
\end{domanda}

\begin{domanda}[Severità ENISA -- \rif{2026-11-02}{02/11}]\label{q:severita}
A breach exposes live location data (behavioural data, DPC~=~2) of fully identifiable users (EI~=~1),
with loss of confidentiality (0.5) and malicious intent (+0.5). Severity?
\begin{opzioni}
  \item 2, medium
  \item 2.5, medium
  \item 3, high
  \item 4, very high
\end{opzioni}
\end{domanda}

\begin{domanda}[Severità ENISA -- \rif{2026-10-12}{12/10}]\label{q:severita-salute}
Medical certificates (DPC~=~4) of employees are sent by mistake to the wrong internal mailing list;
identification is significant (EI~=~0.75); only confidentiality is affected (0.5), no malicious intent.
Severity?
\begin{opzioni}
  \item 2.5, medium
  \item 3, high
  \item 3.5, high
  \item 4.5, very high
\end{opzioni}
\end{domanda}

\begin{domanda}[GDPR, art.~34 -- \rif{2026-11-02}{02/11}]\label{q:minori}
A high-risk breach concerns personal data of children aged 5--12. The Art.~34 communication must be
addressed to:
\begin{opzioni}
  \item the children directly;
  \item the holders of parental responsibility (the parents);
  \item the children's schools;
  \item no one: communication is not required for minors, only notification to the authority.
\end{opzioni}
\end{domanda}

% =====================================================================
\section{ISO/IEC 27001 e risk management}

\begin{domanda}[ISO 27001 -- \rif{2026-10-19}{19/10}]\label{q:annexa}
ISO/IEC 27001:2022 Annex~A contains:
\begin{opzioni}
  \item 114 controls in 14 domains;
  \item 93 controls in 4 themes (organisational, people, physical, technological);
  \item 10 mandatory clauses and no controls;
  \item 34 technological controls only.
\end{opzioni}
\end{domanda}

\begin{domanda}[ISO 27001 -- \rif{2026-10-19}{19/10}]\label{q:certificazione}
About ISO/IEC 27001 certification. \textbf{Select the WRONG statement.}
\begin{opzioni}
  \item Clauses 4--10 are mandatory requirements of the ISMS.
  \item The certificate follows a three-year cycle with surveillance audits.
  \item The certificate guarantees that the organisation will not suffer security incidents.
  \item Certification concerns the organisation and its ISMS, not a single product.
\end{opzioni}
\end{domanda}

\begin{domanda}[Risk management -- \rif{2026-10-26}{26/10}]\label{q:max-cia}
In the course's risk register a risk is scored C~=~1, I~=~1, A~=~3 (scale 1--3). Its level is:
\begin{opzioni}
  \item 1, the minimum;
  \item about 1.7, the average;
  \item 3, the maximum of the three;
  \item 5, the sum.
\end{opzioni}
\end{domanda}

\begin{domanda}[Risk management -- \rif{2026-10-26}{26/10}]\label{q:trattamento}
Which of the following is \textbf{not} a risk treatment option?
\begin{opzioni}
  \item mitigate (modify) the risk with controls;
  \item transfer or share the risk;
  \item avoid the risk by stopping the activity;
  \item ignore the risk without documenting it.
\end{opzioni}
\end{domanda}

\begin{domanda}[Risk management -- \rif{2026-10-26}{26/10}]\label{q:soa}
The Statement of Applicability (SoA):
\begin{opzioni}
  \item lists the Annex~A controls, states whether each is applied and justifies inclusions and
    exclusions;
  \item is the certificate issued by the certification body;
  \item is the GDPR record of processing activities;
  \item is required only for organisations subject to NIS2.
\end{opzioni}
\end{domanda}

\begin{domanda}[ISO vs GDPR -- \rif{2026-10-23}{23/10}]\label{q:impatto}
Which statement about \emph{impact} is \textbf{correct}?
\begin{opzioni}
  \item ISO/IEC 27001 and the ENISA severity methodology measure impact on the same object.
  \item In ISO/IEC 27001 impact is assessed on the organisation; in the GDPR/ENISA approach on the
    rights and freedoms of data subjects.
  \item In ISO/IEC 27001 impact is assessed only on data subjects.
  \item The GDPR does not consider impact.
\end{opzioni}
\end{domanda}

\begin{domanda}[Rischio -- \rif{2026-10-05}{05/10}]\label{q:residuo}
Residual risk is:
\begin{opzioni}
  \item the risk before any control is applied;
  \item the risk remaining after the mitigation measures have been applied;
  \item the risk transferred to an insurer;
  \item the risk that the authority accepts on behalf of the controller.
\end{opzioni}
\end{domanda}

\begin{domanda}[ISO e conformità -- \rif{2026-11-02}{02/11}]\label{q:iso}
A company holds an ISO/IEC 27001 certificate. \textbf{Select the WRONG statement.}
\begin{opzioni}
  \item The certificate can support the demonstration of accountability.
  \item The certificate proves that the company is GDPR and CRA compliant.
  \item The certificate does not exempt the company from data breach notification duties.
  \item Auditors refer to ISO/IEC 27002 to assess how the controls are implemented.
\end{opzioni}
\end{domanda}

% =====================================================================
\section{Cyber Resilience Act}

\begin{domanda}[CRA, ambito -- \rif{2026-10-26}{26/10}]\label{q:cra-ambito}
Which of the following is \textbf{not} within the scope of the CRA?
\begin{opzioni}
  \item a GPS smartwatch for children connected to a cloud backend;
  \item an installable mobile app connected to a remote server;
  \item a medical device certified under the MDR;
  \item a home Wi-Fi router.
\end{opzioni}
\end{domanda}

\begin{domanda}[CRA, transizione -- \rif{2026-11-02}{02/11}]\label{q:cra-unita}
For CRA applicability, the cut-off date of 11 December 2027 is assessed with reference to:
\begin{opzioni}
  \item the date the product model was first launched;
  \item each individual unit placed on the market;
  \item the manufacturer's date of incorporation;
  \item the date of the last firmware update.
\end{opzioni}
\end{domanda}

\begin{domanda}[CRA, segnalazioni -- \rif{2026-11-02}{02/11}]\label{q:cra-final}
Under the CRA, the final report on an actively exploited vulnerability is due:
\begin{opzioni}
  \item within 72 hours of becoming aware of the vulnerability;
  \item within one month of becoming aware of the vulnerability;
  \item no later than 14 days after a corrective or mitigating measure is available;
  \item there is no deadline for the final report.
\end{opzioni}
\end{domanda}

\begin{domanda}[CRA e GDPR -- \rif{2026-11-02}{02/11}]\label{q:cra-gdpr}
The same incident is both a severe incident affecting a product (CRA) and a personal data breach
(GDPR). \textbf{Select the WRONG statement.}
\begin{opzioni}
  \item The CRA early warning is due within 24 hours of awareness.
  \item The GDPR notification is due within 72 hours of awareness.
  \item Sending the CRA early warning to the CSIRT also fulfils the GDPR notification duty.
  \item The two notifications run in parallel, towards different authorities.
\end{opzioni}
\end{domanda}

\begin{domanda}[CRA e AI Act -- \rif{2026-11-02}{02/11}]\label{q:cra-ai}
A product with digital elements that is also a high-risk AI system:
\begin{opzioni}
  \item is outside the CRA, since the AI Act prevails;
  \item is presumed compliant with the AI Act cybersecurity requirements (Art.~15) if it meets CRA
    Annex~I (Parts I and II) and this is shown in the CRA EU declaration of conformity;
  \item must obtain two separate CE markings, one per regulation;
  \item is automatically classified as a critical product under the CRA.
\end{opzioni}
\end{domanda}

\begin{domanda}[Metodo di audit -- \rif{2026-11-06}{06/11}]\label{q:fonti}
In the conformity audit method used in the course, a statement confirmed by two independent and
concordant sources is treated as:
\begin{opzioni}
  \item an assertion, to be taken with caution;
  \item an established fact (evidence);
  \item a finding of non-conformity;
  \item irrelevant, unless signed by the CEO.
\end{opzioni}
\end{domanda}

% =====================================================================
\section{Cybersecurity Act}

\begin{domanda}[ENISA -- \rif{2026-11-09}{09/11}]\label{q:enisa}
Which statement about ENISA is \textbf{correct}?
\begin{opzioni}
  \item ENISA imposes administrative fines on essential entities that fail to report incidents.
  \item ENISA received a permanent mandate with the Cybersecurity Act and supports Member States and EU
    institutions, e.g.\ through the CRA Single Reporting Platform.
  \item ENISA directly issues ISO/IEC 27001 certificates to companies.
  \item ENISA is the national CSIRT of Italy.
\end{opzioni}
\end{domanda}

\begin{domanda}[CSA, certificazione -- \rif{2026-11-09}{09/11}]\label{q:livelli}
The assurance levels of European cybersecurity certificates under the CSA are:
\begin{opzioni}
  \item low, medium, high, very high;
  \item basic, substantial, high;
  \item bronze, silver, gold;
  \item self-assessed, audited, certified.
\end{opzioni}
\end{domanda}

\begin{domanda}[CSA -- \rif{2026-11-09}{09/11}]\label{q:csa-cra}
About the relationship between CSA and CRA. \textbf{Select the WRONG statement.}
\begin{opzioni}
  \item The CSA creates a European framework of certification schemes.
  \item A CSA certificate can help demonstrate conformity with the CRA.
  \item The CSA replaces the CRA for certified products.
  \item EUCC is the first scheme adopted under the CSA and is based on the Common Criteria.
\end{opzioni}
\end{domanda}

% =====================================================================
\section{NIS2, DORA e AI Act}

\begin{nota}
Norme non ancora trattate a lezione (al 11/11/2026): domande basate sull'appello precedente e
sull'inquadramento dei Capitoli~\ref{cap:nis2}--\ref{cap:aiact}.
\end{nota}

\begin{domanda}[NIS2]\label{q:nis2-incidente}
Under NIS2, an \emph{incident} is:
\begin{opzioni}
  \item an event compromising the availability, authenticity, integrity or confidentiality of stored,
    transmitted or processed data or of the services offered by, or accessible via, network and
    information systems;
  \item a breach of security leading to the accidental or unlawful destruction, loss, alteration,
    unauthorised disclosure of, or access to, personal data;
  \item a vulnerability for which there is reliable evidence that a malicious actor has exploited it;
  \item only an event that has actually caused a financial loss above a fixed threshold.
\end{opzioni}
\end{domanda}

\begin{domanda}[DORA]\label{q:dora-strategie}
Digital operational resilience strategy under DORA. \textbf{Select the WRONG statement.}
\begin{opzioni}
  \item It includes the entity's ICT risk tolerance and resilience objectives.
  \item It may include a multi-vendor strategy for ICT third-party service providers.
  \item It includes a communication strategy for ICT-related incidents.
  \item It must be approved by the main ICT third-party service provider before adoption.
\end{opzioni}
\end{domanda}

\begin{domanda}[DORA]\label{q:dora-iniziale}
Under DORA, the initial notification of a major ICT-related incident must be submitted:
\begin{opzioni}
  \item within 72 hours of awareness;
  \item within 4 hours from the classification as major, and no later than 24 hours from awareness;
  \item within 24 hours from the classification as major, with no other limit;
  \item within one month, together with the final report.
\end{opzioni}
\end{domanda}

\begin{domanda}[NIS2 e DORA]\label{q:lex-specialis}
A bank is in a sector covered by both NIS2 and DORA. For ICT incident reporting:
\begin{opzioni}
  \item NIS2 prevails, being a directive;
  \item the bank must report the same incident twice, under both regimes;
  \item DORA applies as lex specialis;
  \item neither applies: the GDPR covers it.
\end{opzioni}
\end{domanda}

% =====================================================================
\section{Tempistiche trasversali}

\begin{domanda}[Tempistiche -- \rif{2026-11-02}{02/11}]\label{q:72h}
About ``72 hours'' in EU law. \textbf{Select the WRONG statement.}
\begin{opzioni}
  \item GDPR: notification to the supervisory authority within 72 hours of awareness of the breach.
  \item NIS2: incident notification within 72 hours of awareness of the significant incident.
  \item DORA: intermediate report within 72 hours of awareness of the incident.
  \item CRA: vulnerability notification within 72 hours of awareness.
\end{opzioni}
\end{domanda}

% =====================================================================
\section{Soluzioni}

\newcommand{\sol}[2]{\item[\ref{#1}.] \textbf{(#2)}.}
\begin{description}[leftmargin=2.5em,labelwidth=2em,font=\normalfont\bfseries]
  \sol{q:regolamento}{b} Il regolamento si applica direttamente; la direttiva vincola sul risultato e va
    recepita (es.\ NIS2 con il D.lgs.~138/2024).
  \sol{q:softlaw}{c} Gli standard ISO sono volontari: diventano rilevanti perché aiutano a dimostrare
    il rispetto della hard law.
  \sol{q:dpa}{b} Trappola segnalata dal docente: nei contratti ``DPA'' è il \emph{Data Processing
    Agreement}; l'autorità è la \emph{data protection authority}.
  \sol{q:anonimi}{c} Togliere nome e cognome non basta e i dati pseudonimizzati restano personali; solo
    i dati davvero anonimi escono dal GDPR.
  \sol{q:gdpr-basi}{c} L'art.~6(1) esclude il legittimo interesse per i trattamenti delle autorità
    pubbliche nell'esecuzione dei loro compiti: per quelli si usa l'interesse pubblico o l'obbligo legale.
  \sol{q:ruoli-fatto}{b} Conta la situazione di fatto, non ciò che dice il contratto: punto su cui il
    docente ha detto che sarà severo.
  \sol{q:autorizzato}{c} Il dipendente è una persona autorizzata che agisce sotto l'autorità del
    titolare; il responsabile è un soggetto esterno.
  \sol{q:legittimo}{b} Il legittimo interesse è del titolare (o di terzi); l'interessato ha diritti
    fondamentali.
  \sol{q:principi}{d} I principi dell'art.~5 sono liceità, correttezza e trasparenza; limitazione della
    finalità; minimizzazione; esattezza; limitazione della conservazione; integrità e riservatezza (più
    l'accountability).
  \sol{q:domestica}{b} È il principio del caso Lindqvist: la pubblicazione su Internet rende i dati
    accessibili a un numero indefinito di persone ed esclude l'esenzione domestica.
  \sol{q:sanzioni}{b} ``Domanda insidiosa'': per sicurezza e data breach il tetto è la metà di quello
    generale (art.~83(4) contro art.~83(5)).
  \sol{q:malattia}{b} Ogni dato relativo alla salute ricade nell'art.~9: nel modello ENISA è DPC~4,
    che si può abbassare solo motivando.
  \sol{q:dpia}{b} Se il rischio residuo resta elevato si consulta preventivamente l'autorità
    (art.~36). Conta il residuo, non l'inerente.
  \sol{q:def-breach}{c} Senza compromissione dei dati personali non c'è violazione: l'email di
    phishing cancellata è al più un evento di sicurezza.
  \sol{q:accesso}{b} La riservatezza è violata con il solo accesso non autorizzato; non presumere
    l'assenza di esfiltrazione se non la si può provare.
  \sol{q:decorrenza}{b} Dalla consapevolezza del titolare, che va però dimostrata: una data di
    consapevolezza tardiva va giustificata con evidenze.
  \sol{q:responsabile}{b} Solo il titolare notifica all'autorità; il responsabile informa il titolare
    senza ingiustificato ritardo (art.~33(2)).
  \sol{q:registro}{b} Il registro documenta \emph{ogni} violazione, anche se non notificata: è la
    prova dell'accountability.
  \sol{q:disponibilita}{b} Anche l'indisponibilità temporanea è un data breach. Se il rischio è basso
    può bastare la registrazione interna.
  \sol{q:formula}{b} È la ``regola d'oro'' da ricordare a memoria.
  \sol{q:severita}{c} $SE = 2 \times 1 + 0{,}5 + 0{,}5 = 3$: alta (è il Gruppo~B del caso Pulse K).
  \sol{q:severita-salute}{c} $SE = 4 \times 0{,}75 + 0{,}5 = 3{,}5$: alta ($3 \le SE < 4$).
  \sol{q:minori}{b} Giuridicamente gli interlocutori sono i titolari della responsabilità genitoriale.
  \sol{q:annexa}{b} 93 controlli in 4 temi (37 organizzativi, 8 persone, 14 fisici, 34 tecnologici).
    La (a) descrive la vecchia versione 2013.
  \sol{q:certificazione}{c} La certificazione attesta il sistema di gestione, non l'assenza di
    incidenti.
  \sol{q:max-cia}{c} Il livello è il massimo di C, I, A: fare la media è un errore tipico.
  \sol{q:trattamento}{d} Accettare un rischio è legittimo, ma va deciso e documentato; ignorarlo no.
  \sol{q:soa}{a} La SoA documenta le decisioni sui controlli ed è il punto in cui si rende conto delle
    esclusioni.
  \sol{q:impatto}{b} Non confondere i due piani: organizzazione (ISO) contro interessati (GDPR/ENISA).
  \sol{q:residuo}{b} Inerente meno mitigazioni: è il residuo che guida la decisione.
  \sol{q:iso}{b} La ISO è uno standard tecnico e non prova la conformità legale: è il suggerimento
    dato da Gianluca per il laboratorio di audit.
  \sol{q:cra-ambito}{c} I dispositivi medici sono soggetti all'MDR, regime verticale che esclude il CRA.
    Le app connesse e commercializzate rientrano: il software è un prodotto.
  \sol{q:cra-unita}{b} Le date di cut-off si valutano per singola unità: unità dello stesso modello
    possono avere obblighi diversi.
  \sol{q:cra-final}{c} Il termine di 14 giorni decorre dalla disponibilità del fix, non dalla
    consapevolezza. Per gli incidenti gravi, invece, il final report è dovuto entro un mese.
  \sol{q:cra-gdpr}{c} Una notifica a un'autorità non adempie gli obblighi verso un'altra: l'early
    warning CRA non è la notifica preliminare GDPR.
  \sol{q:cra-ai}{b} È la presunzione di conformità vista il 02/11, che collega l'art.~15 AI Act e
    l'Allegato~I del CRA.
  \sol{q:fonti}{b} Una fonte è un'asserzione; due fonti concordanti sono un fatto accertato.
  \sol{q:enisa}{b} Le sanzioni NIS2 sono irrogate dalle autorità nazionali; i certificati sono
    rilasciati da organismi di valutazione o autorità nazionali (e ISO non è uno schema UE); il CSIRT
    italiano è presso l'ACN.
  \sol{q:livelli}{b} Base, sostanziale, elevato (\emph{basic, substantial, high}).
  \sol{q:csa-cra}{c} Il CSA è soft law e aiuta a dimostrare la conformità al CRA (hard law), ma non lo
    sostituisce.
  \sol{q:nis2-incidente}{a} (b) è la definizione GDPR di data breach; (c) è la vulnerabilità
    attivamente sfruttata del CRA; (d) è falsa: rileva anche il danno potenziale e nella definizione non
    c'è alcuna soglia.
  \sol{q:dora-strategie}{d} La strategia è definita e approvata dall'organo di gestione dell'entità
    finanziaria; ai fornitori ICT gli obblighi arrivano per via contrattuale.
  \sol{q:dora-iniziale}{b} Doppio limite: 4 ore dalla classificazione come grave e comunque non oltre
    24 ore dalla conoscenza. Le 72 ore di DORA riguardano invece la relazione intermedia.
  \sol{q:lex-specialis}{c} DORA è lex specialis rispetto a NIS2 per le entità finanziarie.
  \sol{q:72h}{c} Le 72 ore di DORA decorrono dalla notifica iniziale, non dalla consapevolezza.
\end{description}
